%% kapa: exam=1 round=24 year=2013 period=2 form=A subject=accounting subject_ko=회계학 seq=12 number=12
%% source: Q-Net 감정평가사 기출문제 (cst00309) artlSeq=5073802 file=extracted/1008/5073802/2013년 제24회 감정평가사 1차시험 기출문제/2교시(A형).hwp
%% conversion: kordoc 4.24.51 -> tools/gen1.py (auto); review=proofread-sonnet; flags=none
\begin{question}{12}
재고자산의 측정에 관한 설명으로 옳지 않은 것은?
\begin{choices}
  \item 표준원가법으로 평가한 결과가 실제원가와 유사하지 않은 경우에는 편의상 표준원가법을 사용할 수 있다.
  \item 개별법은 통상적으로 상호 교환될 수 없는 항목이나 특정 프로젝트별로 생산되고 분리되는 재화 또는 용역에 적용하는 방법이다.
  \item 생물자산에서 수확한 농림어업 수확물로 구성된 재고자산은 순공정가치로 측정하여 수확시점에 최초로 인식한다.
  \item 소매재고법은 이익률이 유사하고 품종변화가 심한 다품종 상품을 취급하는 유통업에서 실무적으로 다른 원가측정법을 사용할 수 없는 경우에 흔히 사용한다.
  \item 후입선출법은 대부분의 경우 실제물량흐름과 반대라는 점, 재고층의 청산 시 수익ㆍ비용 대응구조의 왜곡 등 여러 가지 비판으로 한국채택국제회계기준에서는 인정되지 않고 있다.
\end{choices}
\end{question}
\answer{1}
